\subsection{道路像的拓扑性质}

命题 15. --- 分离拓扑空间中一条道路的像，是紧、可度量化、连通且局部道路连通的拓扑空间。

设 X 为分离拓扑空间，\(c\colon I \to X\) 为连续映射。用 R 表示 I 中的等价关系 \(c(s) = c(t)\)。拓扑空间 \(c(\mathbf{I})\) 是分离的（TG, I, p. 63, 推论 1），空间 I/R 是拟紧的（TG, I, p. 62, 定理 2），因此由 c 得到的双射 I/R \(\to c(\mathbf{I})\) 是同胚（TG, I, p. 63, 推论 2）。于是空间 \(c(\mathbf{I})\) 是紧的、可度量化的（TG, IX, p. 22, 命题 17）、连通的（TG, I, p. 82, 命题 6）且局部道路连通的（III, p. 261, 命题 8）。

定理 1（Hahn 与 Mazurkiewicz）. --- 每个可度量化、紧、非空、连通且局部连通的拓扑空间，都同胚于区间 {[}0,1{]} 的一个商空间。

引理 4. --- 设 K 为紧拓扑空间，R 为由覆盖 K 的 K 的开集组成的集合。存在 K 的一致结构的一个邻偶 V（TG, II, p. 27, 定理 1），使得对每个 \(x \in K\)，\(V(x)\) 都包含于 R 的某个元素之中。

对 K 的每个点 x，存在 K 的一致结构的一个邻偶 \(W_{x}\)，使得 \(\mathrm{W}_{x}(x)\) 包含于 R 的某个元素之中。设 \(V_{x}\) 为 K 的一致结构的满足下列条件的邻偶：\(\overset{2}{\mathrm{V}}_{x}\) 包含于 \(W_{x}\)。诸 \(\mathrm{V}_{x}(x)\) 的内部覆盖 K；由于空间 K 紧，存在 K 的有限子集 F，使得族 \((\mathrm{V}_{y}(y))_{y\in\mathrm{F}}\) 覆盖 K（TG, I, p. 59）。用 V 表示族 \((\mathrm{V}_{y})_{y\in\mathrm{F}}\) 的交；集合 V 是 K 的一致结构的邻偶。对每个点 \(x\in K\)，存在点 \(y\in F\)，使得 x 属于 \(\mathrm{V}_{y}(y)\)。于是集合 \(\mathrm{V}(x)\) 包含于 \(\overset{2}{\mathrm{V}}_{y}(y)\)，从而包含于 R 的某个元素之中。

引理 5. --- 设 X 与 Y 为一致空间，f 为从 X 到 Y 中的映射。设 F 为由 X 的闭子集组成的、覆盖 X 且具有下列性质的集族：

\begin{enumerate}
\def\labelenumi{(\roman{enumi})}
\item
  存在一个集合 \(F_{0} \in F\)，它与所有集合 \(F \in F\) 相交；
\item
  对 X 的一致结构的每个邻偶 U，不是 U 阶小集的集合 \(F \in F\) 只有有限多个；
\item
  对 Y 的一致结构的每个邻偶 V，集合 \(F \in F\)（使得 \(f(F)\) 不是 V 阶小集的）只有有限多个。
\end{enumerate}

那么，若 f 在每个集合 \(F \in F\) 上的限制都连续，则 f 连续。

设 x 为 X 的一点。我们来证明 f 在 x 处连续。存在集合 \(F_{1} \in F\)，使得 \(x \in F_{1}\)。\(f\) 在 \(F_{0} \cup F_{1}\) 上的限制是连续的（TG, I, p. 19, 命题 4）。把 \(F_{0}\) 换成 \(F_{0} \cup F_{1}\)，可化归为 \(x \in F_{0}\) 的情形。

设 V 为 Y 的一致结构的一个邻偶。选取该一致结构的一个邻偶 \(V'\)，使得 \(V' \subset V\)。由于 f 在 \(F_0\) 上的限制连续，存在 X 的一致结构的一个邻偶 U，使得 \(f(z) \in \mathrm{V}'(f(x))\)（对所有 \(z \in \mathrm{F}_0 \cap \mathrm{U}(x)\)）。设 \(U'\) 为 X 的一致结构的满足 \(U' \subset U\) 的一个邻偶。用 A 表示 \(F_0\)、诸集合 \(F \in F\) 中不是 \(U'\) 阶小集者、以及诸使 \(f(\mathrm{F})\) 不是 \(V'\) 阶小集的集合的并。按假设，A 是属于 F 的有限多个集合的并，而 f 在 A 上的限制连续（同前所引）。因此存在一个邻域 W（x 的，在 X 中），它包含于 \(U'(x)\)，且 \(f(y) \in \mathrm{V}(f(x))\)（对 \(y \in A \cap W\)）。为完成证明，只需证明也有 \(f(y) \in \mathrm{V}(f(x))\)（对每个点 \(y \in (\mathrm{X} - \mathrm{A}) \cap \mathrm{W}\)）。设 y 为这样一个点。设 F 为 F 的满足 \(y \in F\) 的一个元素。由 A 的定义，F 是 \(U'\) 阶小集，且 \(f(\mathrm{F})\) 是 \(V'\) 阶小集。按假设，F 与 \(F_0\) 相交。设 \(z \in F \cap F_0\)。我们有 \(z \in U'(y)\)，因为 F 是 \(U'\) 阶小集；又有 \(y \in U'(x)\)，因为 W 包含于 \(U'(x)\)；于是 \(z \in U'(x)\)，更有 \(z \in U(x)\)。但这时，由于 z 属于 \(F_0\)，我们有 \(f(z) \in V'(f(x))\)。此外 \(f(\mathrm{F})\) 是 \(V'\) 阶小集，于是 \(f(y) \in V'(f(z))\)。由此得出 \(f(y) \in V'(f(x))\)，最终有 \(f(y) \in V(f(x))\)。引理 5 证毕。

现在来证明定理 1。设 X 为非空的紧、连通且局部连通的度量空间。这样的空间是道路连通且局部道路连通的（III, p. 267, 推论 2）。由 III, p. 270 的推论与例，存在从康托尔三分集 K（TG, IV, p. 9, 例）到 X 上的连续满射 f。我们将构造一个连续开拓 g（f 的，到 {[}0,1{]} 上），这就证明了定理 1。

设 \(\varepsilon\) 为实数 \(>0\)。X 的直径 \(\leqslant \varepsilon\) 的开的道路连通子集覆盖 X。用 \(\mathcal{R}\) 表示它们在 \(f\) 下的原像组成的集合；这是由覆盖 K 的 K 的开子集组成的集合。由 III, p. 272 的引理 4，存在实数 \(\alpha > 0\)，使得 K 中半径为 \(\alpha\) 的每个闭球都包含于 \(\mathcal{R}\) 的一个元素之中。特别地，若 \(t\) 与 \(t'\) 为 K 中满足 \(|t - t'| \leqslant \alpha\) 的点，则存在 X 中连接 \(f(t)\) 与 \(f(t')\) 的道路，其像的直径 \(\leqslant \varepsilon\)。

前一段使我们可以归纳地构造严格递增序列 \((n_{k})_{k\geqslant0}\)，其项为整数 \(\geqslant0\)，且具有下列性质：对每个整数 \(k\geqslant0\) 与每对点 \((t,t')\)（K 的、满足 \(|t-t'| \leqslant3^{-n_{k}}\) 的），存在 X 中连接 \(f(t)\) 与 \(f(t')\) 的道路，其像的直径 \(\leqslant2^{-k}\)。K 在 {[}0,1{]} 中的余集是两两不交的开区间组成的族 \((\mathrm{I}_{n,p})\) 的并，其中 n 取遍整数 \(\geqslant0\) 组成的集合，而 p 取遍 1 与 \(2^{n}\) 之间的整数组成的集合（TG, IV, p. 9, 例）。考虑这些区间 \(I_{n,p}\) 之一，并把它写成 {]}a,b{[}。点 a 与 b 属于 K，且有 \(b-a=3^{-n-1}\)。如下定义区间 \(I_{n,p}\) 上的函数 g：选取 X 中连接 \(f(a)\) 与 \(f(b)\) 的道路 c，其像的直径 \(\leqslant2^{-k}\)（当 \(n_{k}\leqslant n<n_{k+1}\) 时），并令 \(g(t)=c(\frac{t-a}{b-a})\)（\(t\in I_{n,p}\)）。这样定义的函数 \(g\colon[0,1]\to X\) 是 f 的开拓。它在 K 上以及在每个闭区间 \(\overline{I_{n,p}}\) 上都连续。这些闭区间与 K 相交。此外，对每个实数 \(\varepsilon>0\)，区间 \(\overline{I_{n,p}}\) 中长度 \(>\varepsilon\) 的只有有限多个，而区间 \(\overline{I_{n,p}}\) 中在 g 下的像的直径 \(>\varepsilon\) 的也只有有限多个。由引理 5，映射 g 连续。
% semantic-fragments: legacy-input-seam
\relax{}
